\section{Comparison to Enacted Maps}
\label{sec:comparison}

We compare algorithmic maps to enacted 2020-cycle state house maps for the
38 states that had enacted their new maps by the time of this analysis.
The comparison focuses on three metrics: compactness (Polsby-Popper),
population balance, and county splits.

\subsection{Compactness Comparison}
\label{sec:comparison:compactness}

The algorithmic maps are more compact (higher mean PP) than enacted maps in
34 of 38 states (89\%). The mean advantage is $+0.081$ PP points
(algorithmic $\overline{PP} = 0.381$ vs.\ enacted $\overline{PP} = 0.300$
for the same 38 states).

The four states where enacted maps are more compact than algorithmic maps
(Vermont, Hawaii, Alaska, and Wyoming) are all small, geographically compact
states where the enacted maps are essentially forced to be compact by the
state's shape and population distribution. In these states, the enacted
map and the algorithmic map are nearly identical, and the small PP differences
are within the noise of the metric.

\subsection{Population Balance Comparison}
\label{sec:comparison:balance}

All 50 algorithmic maps achieve $\leq 0.5\%$ maximum deviation. Of the
38 enacted maps, 26 (68\%) achieve $\leq 0.5\%$; 12 (32\%) have maximum
deviations in the $0.5$--$5.0\%$ range permitted by \emph{Reynolds v.\ Sims}.
No enacted map exceeds $5\%$.

The algorithmic maps are thus more population-balanced than enacted maps
by a statistically significant margin: a paired t-test of maximum deviations
gives $t = 3.8$, $p < 0.001$.

\subsection{County Splits Comparison}
\label{sec:comparison:splits}

As reported in Section~\ref{sec:results:splits}, the algorithmic maps split
counties fewer times than enacted maps in 30 of 38 states (79\%).
The mean reduction of 18\% is statistically significant ($t = 4.2$, $p < 0.001$).

The magnitude of the reduction varies dramatically across states.
States with historically aggressive gerrymandering (Pennsylvania, Wisconsin,
Ohio, North Carolina) show the largest reductions (30--41\%), while states
with citizen redistricting commissions (California, Colorado, Michigan) show
smaller reductions (5--12\%). This pattern is consistent with the hypothesis
that enacted maps in gerrymander-prone states deliberately crack communities
across county lines, and that the algorithmic maps---which have no partisan
inputs---avoid such cracking.

\subsection{National Summary}
\label{sec:comparison:summary}

\begin{table}[h]
\centering\small
\begin{tabular}{lrr}
\toprule
Metric & Algorithmic & Enacted \\
\midrule
Mean PP (38 states) & 0.381 & 0.300 \\
States more compact than enacted & 34/38 & --- \\
Max population deviation $\leq 0.5\%$ & 50/50 & 26/38 \\
Mean county splits & 22.4 & 27.4 \\
States with fewer splits than enacted & 30/38 & --- \\
\bottomrule
\end{tabular}
\caption{Aggregate comparison of algorithmic vs.\ enacted state house maps
  (38 states with available enacted 2020-cycle maps). The algorithmic maps
  are more compact, more population-balanced, and split fewer counties in
  the large majority of states.}
\label{tab:comparison-summary}
\end{table}
